\section{Limitations and Future Directions}
\label{sec:future}

While our results demonstrate the promise of hybrid physics-ML approaches for ENSO forecasting, several limitations of the current work suggest directions for future research.

\subsection{Data and Evaluation Limitations}

Our evaluation relies on a single train-test split (1979--2001 for training, 2002--2022 for testing), which provides a rigorous out-of-sample assessment but does not characterize variability across different climate regimes. The training period includes the major 1982--83 and 1997--98 El Ni\~no events, while the test period includes the 2015--16 event and extended La Ni\~na conditions of 2020--2022. Model rankings might differ with alternative splits that allocate these events differently.

Rolling-window cross-validation would provide more robust model selection by averaging performance across multiple train-test splits. However, the relatively short observational record (44 years) limits the number of independent folds possible while maintaining sufficient training data. Extending the analysis to earlier periods using reconstructed indices would expand the available data but introduce additional uncertainty from the reconstruction methodology.

The current framework uses only ORAS5~\cite{zuo2019oras5} reanalysis for training and evaluation. Multiple reanalysis products (ERA5~\cite{hersbach2020era5}, GODAS, SODA) provide alternative representations of the historical climate state, each with distinct biases and uncertainty characteristics. Training on multiple products or evaluating across products would assess robustness to observational uncertainty. Our codebase supports multi-dataset evaluation, but systematic comparison remains for future work.

\subsection{Architectural Extensions}

The transformer failure in our experiments does not preclude success with modified architectures or training procedures. Recent work on reservoir computing for El Ni\~no prediction~\cite{guardamagna2025reservoir} demonstrates that alternative recurrent architectures can achieve high skill, suggesting that the key limitation is not recurrence per se but the combination of architecture, data volume, and inductive biases. Several remediation strategies merit investigation.

Pre-training on large climate model ensembles before fine-tuning on observations could provide transformers with the data volume they require. CMIP6~\cite{eyring2016cmip6} simulations span thousands of model-years across 100+ models, offering abundant data with realistic (though not identical) ENSO dynamics. The challenge lies in the domain shift between simulated and observed climate, which our residual model experiments suggest can be detrimental for flexible architectures.

Physics-informed attention mechanisms could encode domain knowledge as architectural constraints. Rather than allowing unrestricted self-attention, the attention weights could be masked or regularized to follow known teleconnection patterns, similar to our NXRO-Attentive model but applied within a full transformer architecture.

Hybrid transformer architectures that retain the physics-informed linear operator while using transformer layers only for the residual might combine the benefits of both approaches. The linear operator would provide the inductive bias needed for small-data learning, while the transformer would offer flexible capacity for capturing complex residual dynamics.

For graph-based models, allowing the adjacency matrix to be learned during training rather than fixed from XRO could discover novel teleconnection patterns, similar to data-driven discovery of coordinates and governing equations~\cite{champion2019data}. Sparsity-inducing regularization would be necessary to prevent the learned graph from becoming fully connected. Multi-scale graph architectures with hierarchical structure could capture both local interactions and long-range teleconnections. Additionally, Neural Flows~\cite{bilos2021neural} offer a computationally efficient alternative to Neural ODEs that could enable scaling to higher-dimensional state spaces.

\subsection{Stochastic Modeling Extensions}

The Stage 3 failure reveals that climate model ensemble spread does not directly transfer to neural forecast uncertainty, but modifications might recover its value. Scaling the simulation-derived noise by a learned factor ($\sigma_{\text{scaled}} = \alpha \cdot \sigma_{\text{simulations}}$ with $\alpha \ll 1$) could reduce the over-dispersion while retaining information about structural uncertainty. The scaling factor could be optimized via the same likelihood objective used in Stage 2.

Selecting a subset of high-quality simulations rather than using all 100 CMIP6 members might improve Stage 3 performance. Ensemble members with smaller biases relative to observations would contribute noise estimates more relevant for well-tuned neural models. Alternatively, using only ensemble reanalyses (such as ERA5-ens or ORAS5-ens) that share the observational constraints of the training data could reduce the distribution mismatch.

Hybrid noise models that combine residual-based and simulation-based uncertainty could capture both aleatoric uncertainty (from atmospheric noise, represented by model residuals) and epistemic uncertainty (from structural model limitations, represented by simulation spread). The relative weighting could vary with lead time, with residual-based uncertainty dominating at short leads and simulation-based uncertainty becoming more relevant at long leads where model structural errors accumulate.

Deep ensembles~\cite{lakshminarayanan2017simple}---training multiple NXRO models with different random initializations---provide an alternative approach to epistemic uncertainty quantification that does not rely on external simulations. The spread across ensemble members captures uncertainty arising from the non-convex optimization landscape and finite training data.

\subsection{Interpretability and Scientific Discovery}

The learned components of NXRO models encode information about climate dynamics that could yield scientific insights beyond forecast improvement.

The attention weights in NXRO-Attentive reveal which climate modes the model considers relevant for predicting each variable at each time step. Visualizing these attention patterns across different ENSO phases could discover state-dependent teleconnection relationships not captured by the static linear operator. Comparisons with physical understanding of ENSO dynamics would validate whether the learned patterns correspond to known mechanisms or suggest novel hypotheses.

The graph convolution weights in NXRO-Graph provide interpretable coupling strengths between climate indices. Comparing these learned weights to the XRO-derived adjacency and to physical coupling estimates from recharge oscillator theory could identify discrepancies that reveal limitations of current physical understanding.

The MLP residual in NXRO-Res captures dynamics that the linear operator misses. Analyzing the residual's dependence on state amplitude, season, and climate mode configuration could identify the functional forms of missing nonlinearities. If the residual approximates specific polynomial or other structured functions, this would suggest extensions to the XRO basis.

\subsection{Operational Considerations}

Deploying NXRO models for operational ENSO forecasting requires addressing practical considerations beyond forecast skill.

Real-time forecasting pipelines must handle data latency, quality control, and automated retraining as new observations become available. The computational efficiency of NXRO models (training in minutes on standard hardware, inference in milliseconds) makes operational deployment feasible, but robust software infrastructure is necessary.

Communicating probabilistic forecasts to decision-makers requires careful attention to visualization and uncertainty representation. The ensemble plumes shown in our results provide one approach, but alternative representations (probability of exceedance curves, categorical probability forecasts) may better serve specific applications in agriculture, water management, or disaster preparedness.

Comparison with operational forecast centers (IRI, CPC, ECMWF, JMA) would contextualize NXRO performance relative to state-of-the-art dynamical and statistical systems. Such comparisons require real-time forecast submission and evaluation over multiple years to assess skill in operational conditions rather than retrospective analysis.

\subsection{Broader Applications}

The NXRO framework extends naturally to other climate forecasting problems where physics-based models exist but have known limitations.

The Indian Ocean Dipole~\cite{saji1999dipole} exhibits recharge-oscillator dynamics similar to ENSO, with thermocline depth anomalies in the eastern Indian Ocean driving SST variability. Adapting NXRO to IOD prediction would require redefining the state variables and potentially the RO nonlinear basis, but the core framework of physics-informed linear dynamics plus neural residual would apply directly.

The Atlantic Multi-decadal Oscillation operates on longer timescales (60--80 years), presenting challenges for training and evaluation given the limited observational record. Transfer learning from climate model simulations might be more valuable for AMO than for ENSO, as the models provide the only source of multiple complete cycles.

Sub-seasonal to seasonal (S2S) prediction~\cite{vitart2017subseasonal} of temperature and precipitation anomalies represents a broader application domain where ENSO serves as a source of predictability. Extending NXRO to predict regional climate impacts conditioned on ENSO state would connect the framework to applications with direct societal relevance.

Finally, the hybrid physics-ML approach embodied by NXRO has potential beyond climate science. Any domain with well-understood but imperfect physics---fluid dynamics, molecular simulation, epidemiology---could benefit from architectures that encode physical structure while learning corrections from data~\cite{daw2022physics,beucler2021enforcing}. The principles we have identified---physics provides sample-efficient inductive bias, ML provides flexibility, end-to-end optimization aligns training with objectives---appear general rather than domain-specific, as evidenced by the success of similar hybrid approaches across scientific domains~\cite{willard2022integrating,yin2021augmenting}.
